%\documentclass[handout,xcolor=pdftex,dvipsnames,table,mathserif]{beamer}
\documentclass[xcolor=pdftex,dvipsnames,table,mathserif]{beamer}
\usepackage{subfigure}
\usepackage{amsbsy}
\usepackage{tikz}
\usetikzlibrary{arrows}
\usepackage{amsmath,graphicx,dsfont,color}
\usepackage{amsfonts}
\usepackage{amssymb}
\usepackage{array}

\bibliographystyle{apalike}

\setbeamertemplate{bibliography item}{\insertbiblabel}
\setbeamertemplate{bibliography entry title}{}
\setbeamertemplate{bibliography entry location}{}
\setbeamertemplate{bibliography entry note}{}

\DeclareMathOperator*{\argmin}{arg\,min}
\DeclareMathOperator*{\argmax}{arg\,max}
%Definitiona

\newcommand{\x}{\mathbf{x}}
\newcommand{\X}{\mathbf{X}}
\newcommand{\W}{\mathbf{W}} %Weight
\newcommand{\bais}{\mathbf{b}}%Bais
\newcommand{\act}{\texttt{g}}%Activation
\newcommand{\loss}{L}
\newcommand{\pdata}{\hat{p}_{\texttt{data}}}
\newcommand{\nsize}{n}
\newcommand{\param}{\boldsymbol{\theta}}
\newcommand{\featmap}{\boldsymbol{\phi}}
\newcommand{\EV}{\mathbb{E}}







\usepackage{physics}

\graphicspath{{../graphics/}}

\AtBeginSection[]{
  \begin{frame}{Contents}
    \tableofcontents[currentsection, hideothersubsections]
  \end{frame}
}

\AtBeginSubsection[]{
  \begin{frame}{Contents}
    \tableofcontents[currentsection, subsectionstyle=show/shaded/hide]
  \end{frame}
}

\setbeamertemplate{footline}[frame number]{}
\setbeamertemplate{navigation symbols}{}
\setbeamertemplate{section in toc}[square]
\setbeamertemplate{items}[square]

%% For image credits on image bottom right
\usepackage[absolute,overlay]{textpos}
\setbeamercolor{framesource}{fg=gray}
\setbeamerfont{framesource}{size=\tiny}
\newcommand{\source}[1]{\begin{textblock*}{4cm}(8.7cm,8.6cm)
    \begin{beamercolorbox}[ht=0.5cm,right]{framesource}
      \usebeamerfont{framesource}\usebeamercolor[fg]{framesource} Credits: {#1}
    \end{beamercolorbox}
\end{textblock*}}

\title{L'apprentissage profond en image biomédicale}
\author{E. Decencière}
\date{MINES ParisTech\\
  PSL Research University\\
  Center for Mathematical Morphology
}
\titlegraphic{\includegraphics[height=1.7cm]{logoemp}}

\useinnertheme{rounded}
\usecolortheme{rose}

%%%%%%%%%%%%%%%%%%%%%%%%%%%%%%%%%%%%%%%%%%%%%%%%%%
%%%%%%%%%%%%%%%%%%%%%%%%%%%%%%%%%%%%%%%%%%%%%%%%%%
\begin{document}
\begin{frame}
  \titlepage
\end{frame}


%%%%%%%%%%%%%%%%%%%%%%%%%%%%%%%%%%%%
\section{Introduction}
%%%%%%%%%%%%%%%%%%%%%%%%%%%%%%%%%%%%

%%%%%%%%%%%%%%%%%%%%%%%%%%%%%%%%%%%%
\begin{frame}{Imagerie biomédicale}

  \begin{block}{Définition}
    Méthodes utilisées en médecine et biologie pour obtenir des images.
  \end{block}

Nous n'aborderons pas les autres données biomédicales, comme les électroencéphalogrammes, etc.

\end{frame}


%%%%%%%%%%%%%%%%%%%%%%%%%%%%%%%%%%%%
\begin{frame}{Différentes modalités pour différentes applications}

\begin{columns}
  \begin{column}{.5\textwidth}
\begin{itemize}[<+->]
\item Enormes progrès en un siècle
\item Nombreuses modalités, basées sur des phénomènes physiques différents:

  \begin{itemize}
  \item  lumière visible
  \item
    rayons X
  \item
    magnétisme
  \item
    physique nucléaire,
  \item
    ondes sonores

  \end{itemize}


\item Nombre et taille des examens en augmentation
\end{itemize}

  \end{column}


  \begin{column}{.5\textwidth}
\begin{figure}[ht]
  \centering
  \includegraphics[width=0.5\textwidth]{first_medical_xray}
  \caption*{1895: Wilhelm Röntgen  radiographie la main de sa femme.}
    \source{Wikipedia.org}
\end{figure}

  \end{column}
\end{columns}


\end{frame}


%%%%%%%%%%%%%%%%%%%%%%%%%%%%%%%%%%%%
\begin{frame}{Rayons X}


  \begin{figure}[ht]
    \centering
    \includegraphics[width=0.7\textwidth]{ct.jpg}\\
    Scanner X
    \source{Wikipedia.org}
  \end{figure}

\end{frame}


%%%%%%%%%%%%%%%%%%%%%%%%%%%%%%%%%%%%
\begin{frame}{Résonance magnétique nucléaire}

  \begin{figure}[ht]
    \centering
    \includegraphics[width=0.6\textwidth]{irm.jpg}\\
    IRM
    \source{Wikipedia.org}
  \end{figure}

\end{frame}

%%%%%%%%%%%%%%%%%%%%%%%%%%%%%%%%%%%%
\begin{frame}{Microscopie optique}

  \begin{figure}[ht]
    \centering
    \includegraphics[width=0.7\textwidth]{histo.jpg}\\
    Histologie
    \source{Wikipedia.org}
  \end{figure}


\end{frame}

%%%%%%%%%%%%%%%%%%%%%%%%%%%%%%%%%%%%
\begin{frame}{Rétinographie}

\begin{columns}
  \begin{column}{.5\textwidth}
    \begin{figure}[ht]
      \centering
      \includegraphics[width=\textwidth]{fundus_camera.jpg}\\
Rétinographe
    \end{figure}

  \end{column}

  \begin{column}{.5\textwidth}
    \begin{figure}[ht]
      \centering
      \includegraphics[width=\textwidth]{fundus.jpg}\\
Rétine\source{wikipedia.org}
    \end{figure}

  \end{column}
\end{columns}

\end{frame}



%%%%%%%%%%%%%%%%%%%%%%%%%%%%%%%%%%%%
\begin{frame}{Ultra-sons}

  \begin{figure}[ht]
    \centering
    \includegraphics[width=0.7\textwidth]{us.jpg}\\
Echographie
    \source{Wikipedia.org}
  \end{figure}


\end{frame}


%%%%%%%%%%%%%%%%%%%%%%%%%%%%%%%%%%%%%%%%%%%%%%%%%%
\section{Défis et difficultés}


%%%%%%%%%%%%%%%%%%%%%%%%%%%%%%%%%%%%
\begin{frame}{Défis}

  \begin{itemize}[<+->]
  \item Croissance du nombre d'examens
  \item Croissance de leur taille
  \item Stabilité ou décroissance du nombre de médecins
  \end{itemize}

\pause

\begin{block}{}
  On va vers une automatisation croissante de l'analyse des images médicales.
\end{block}

\end{frame}



%%%%%%%%%%%%%%%%%%%%%%%%%%%%%%%%%%%%
\begin{frame}{Difficultés pour l'analyse automatique des images biomédicales}

  \begin{itemize}[<+->]
  \item Protection de la vie privée des patients
  \item Responsabilité
  \item Variabilité de l'état sain et des états pathologiques
  \end{itemize}

\pause

  \begin{block}{}
    Aujourd'hui on s'oriente surtout vers une \alert{aide} à l'analyse.
  \end{block}

\end{frame}

%%%%%%%%%%%%%%%%%%%%%%%%%%%%%%%%%%%%
\begin{frame}{Dépistage de la rétinopathie diabétique}

Résultats du projet ANR-TeleOphta~\cite{decenciere_teleophta_2013}.

  \begin{columns}
    \begin{column}{.5\textwidth}
  \begin{figure}[ht]
    \centering
    \includegraphics[width=\textwidth]{fundus1}
    \caption*{Image normale}
  \end{figure}

    \end{column}

    \begin{column}{.5\textwidth}

\begin{figure}[ht]
  \centering
  \includegraphics[width=\textwidth]{ret-dr}\\
  Rétine avec signes de rétinopathie diabétique
\end{figure}


    \end{column}
  \end{columns}

      \source{OPHDIAT database}

\end{frame}


%%%%%%%%%%%%%%%%%%%%%%%%%%%%%%%%%%%%
\begin{frame}{Dépistage de la rétinopathie diabétique}

\begin{figure}[ht]
  \centering
  \includegraphics[width=0.7\textwidth]{ret-oedeme}
  \caption*{Oedème papillaire}
\end{figure}

      \source{OPHDIAT database}

\end{frame}


%%%%%%%%%%%%%%%%%%%%%%%%%%%%%%%%%%%%%%%%%%%%%%%%%%
\section{Apprentissage profond}

%%%%%%%%%%%%%%%%%%%%%%%%%%%%%%%%%%%%
\begin{frame}{Applications}

\begin{itemize}[<+->]
\item Image $\longmapsto$ Grandeur
\item Image $\longmapsto$ Catégorie
\item Image $\longmapsto$ Partition (segmentation)
\item Image $\longmapsto$ Image
\item ? $\longmapsto$ Images
\end{itemize}

\end{frame}

%%%%%%%%%%%%%%%%%%%%%%%%%%%%%%%%%%%%
\begin{frame}{Image $\longmapsto$ Grandeur}

\begin{block}{Exemple: degré de sévérité de la rétinopathie diabétique}
De 0 (pas de signes) à 5.
\end{block}


\begin{figure}[ht]
  \centering
  \includegraphics[width=0.7\textwidth]{ret-dr}\\
  Rétine avec signes de rétinopathie diabétique
\end{figure}
      \source{OPHDIAT database}

\end{frame}


%%%%%%%%%%%%%%%%%%%%%%%%%%%%%%%%%%%%
\begin{frame}{Image $\longmapsto$ Catégorie}

  \begin{figure}[ht]
  \centering
  \includegraphics[width=\textwidth]{pcam.jpg}
\end{figure}

\begin{itemize}
\item Catégorie: tissu tumoral ou pas.
\item Remarque: en histologie, les images sont très souvent gigantesques (plusieurs milliards de pixels)
\end{itemize}

\end{frame}


%%%%%%%%%%%%%%%%%%%%%%%%%%%%%%%%%%%%
\begin{frame}{Image $\longmapsto$ Partition (segmentation)}

Segmentation d'images histologiques de la peau.

\end{frame}


%%%%%%%%%%%%%%%%%%%%%%%%%%%%%%%%%%%%
\begin{frame}{? $\longmapsto$ Images}

Décodage d'un signal cérébral (IRM fonctionnel) pour reconstruire l'image qui est vue par le sujet~\cite{shen_deep_2019}.

\begin{figure}[ht]
  \centering
  \includegraphics[width=\textwidth]{im_from_fMRI}
\end{figure}


\end{frame}



%%%%%%%%%%%%%%%%%%%%%%%%%%%%%%%%%%%%%%%%%%%%%%%%%%
\section{Conclusion}

%%%%%%%%%%%%%%%%%%%%%%%%%%%%%%%%%%%%
\begin{frame}{Promesses et premiers résultats}

\begin{itemize}[<+->]
\item Sur de nombreuses bases de données de \textit{recherche}, l'apprentissage profond dépasse les performances des spécialistes.
\item Le passage à la clinique pose souvent des problèmes.
\item De nombreuses compagnies travaillent sur le sujet.
\item Des systèmes ont déjà été certifiés.
\end{itemize}


\end{frame}

%%%%%%%%%%%%%%%%%%%%%%%%%%%%%%%%%%%%
\begin{frame}{Obstacles}

\begin{itemize}[<+->]
\item Données!
  \begin{itemize}
  \item Constituer des bases suffisamment grandes ou s'en passer?
  \item Annotation
  \end{itemize}
  \item Interprétabilité
\end{itemize}

\end{frame}


%%%%%%%%%%%%%%%%%%%%%%%%%%%%%%%%%%%%
\begin{frame}{Dangers}

  \begin{itemize}[<+->]
  \item Biais
  \item Manque d'interprétabilité
  \item Vie privée?
  \end{itemize}

\end{frame}




%%%%%%%%%%%%%%%%%%%%%%%%%%%%%%%%%%%%%%%%%%%%%%%%%%
\section*{References}

%%%%%%%%%%%%%%%%%%%%%%%%%%%%%%%%%%%%%%%%%%%%%%%%%%

\frame[allowframebreaks]{

  \scriptsize

  \frametitle{References}


  %\bibliographystyle{amsalpha}
  %\bibliographystyle{apalike}

  \bibliography{../../edf.bib}

  \normalsize

}

%%%%%%%%%%%%%%%%%%%%%%%%%%%%%%%%%%%%%%%%%%%%%%%%%%
%%%%%%%%%%%%%%%%%%%%%%%%%%%%%%%%%%%%%%%%%%%%%%%%%%
%%%%%%%%%%%%%%%%%%%%%%%%%%%%%%%%%%%%%%%%%%%%%%%%%%

\end{document}
